\chapter{Your Python Toolkit and Environment Setup}
\label{ch:python}

\chapterguide%
  {Create folders and run basic \textbf{terminal commands}.}
  {organize data and notebooks, configure an environment, install dependencies, and run labs reproducibly.}

\section{The course project layout}

Use \texttt{UTP-MACHINE-LEARNING/} as the project root, with these folders:

\noindent
\begin{minipage}[t]{0.49\linewidth}
\begin{lstlisting}[style=mlterminal]
all_datasets/
  ames_housing_dataset.csv
  bike_sharing_dataset_day.csv
  bike_sharing_dataset_hour.csv
  breast_cancer_dataset.csv
  credit_card_fraud_dataset.csv
  heart_failure_dataset.csv
  mall_customers_dataset.csv
  medical_cost_personal_dataset.csv
  optical_recognition_handwritten_digits.zip
  red_wine_quality_dataset.csv
  student_performance_dataset.csv
  titanic_dataset.csv
\end{lstlisting}
\end{minipage}\hfill
\begin{minipage}[t]{0.49\linewidth}
\begin{lstlisting}[style=mlterminal]
all_experiments/
  lab01_features_worked.ipynb
  lab02_supervised_workflow_worked.ipynb
  lab03_supervised_models_worked.ipynb
  lab04_clustering_pca_worked.ipynb
  lab05_anomaly_limited_labels_worked.ipynb
  lab06_reinforcement_worked.ipynb
  lab07_deep_learning_cnn_worked.ipynb
  lab08_deep_learning_lstm_worked.ipynb
  predict_passenger.py
models/
results/lab01/ ... lab08/
requirements.txt
\end{lstlisting}
\end{minipage}

{\footnotesize Keep the digits ZIP intact. Each lab has one worked notebook; \texttt{predict\_passenger.py} uses the saved Lab~2 model. Notebooks write models and cards to \texttt{models/} and evidence to \texttt{results/}; do not edit generated outputs.\par}

\section{Download the datasets}

Click a \textbf{dataset name} to download it, then extract the files into \textbf{\texttt{all\_datasets/}}.

\begin{mltable}
\footnotesize
\renewcommand{\arraystretch}{1.0}
\setlength{\tabcolsep}{4pt}
\begin{tabularx}{\linewidth}{@{}>{\raggedright\arraybackslash}p{0.40\linewidth}>{\raggedright\arraybackslash}X@{}}
\toprule
\rowcolor{MLPaleBlue!92}
\textbf{Kaggle Dataset} & \textbf{Description} \\
\midrule
\href{https://www.kaggle.com/datasets/uciml/breast-cancer-wisconsin-data}{Breast Cancer Wisconsin} & Lab~1: scaling; Lab~2: thresholds; Lab~3: KNN, naive Bayes, SVM; Lab~5: limited labels \\
\href{https://www.kaggle.com/datasets/mirichoi0218/insurance}{Medical Cost Personal Data} & Lab~1: pipelines; Lab~2: linear and ridge regression \\
\href{https://www.kaggle.com/datasets/larsen0966/student-performance-data-set}{Student Performance} & Lab~1: data quality, prediction timing, feature selection \\
\href{https://www.kaggle.com/datasets/yasserh/titanic-dataset}{Titanic Dataset} & Lab~1: feature engineering; Lab~2: baselines, tuning, saving; Lab~3: trees \\
\href{https://www.kaggle.com/datasets/lakshmi25npathi/bike-sharing-dataset}{Bike Sharing} & Lab~2: daily/hourly time splits; Lab~8: LSTM forecasting \\
\href{https://www.kaggle.com/datasets/fedesoriano/heart-failure-prediction}{Heart Failure Prediction} & Lab~2 Try-it: mixed-type classification \\
\href{https://www.kaggle.com/datasets/shashanknecrothapa/ames-housing-dataset}{Ames Housing} & Lab~3: random forests and gradient boosting \\
\href{https://www.kaggle.com/datasets/uciml/red-wine-quality-cortez-et-al-2009}{Red Wine Quality} & Lab~3: ensembles and feature importance \\
\href{https://www.kaggle.com/datasets/vjchoudhary7/customer-segmentation-tutorial-in-python}{Mall Customers} & Lab~4: K-means and DBSCAN \\
\href{https://www.kaggle.com/datasets/yadavhim/optical-recognition-of-handwritten-digits}{Optical Recognition of Handwritten Digits} & Lab~4: PCA; Lab~7: CNNs \\
\href{https://www.kaggle.com/datasets/miadul/credit-card-fraud-detection-dataset}{Credit Card Fraud Detection} & Lab~5: anomaly detection and imbalanced evaluation \\
\bottomrule
\end{tabularx}
\end{mltable}

\subsection{Download steps}

\begin{mlsteps}
  \item Open the link, click \textbf{Download}, and \textbf{extract} the archive.
  \item Copy the required \textbf{CSV files or images} into \textbf{\texttt{all\_datasets/}}.
  \item Check that your notebook's \textbf{file paths} match the filenames, including \textbf{capitalization}.
\end{mlsteps}

\begin{pitfall}
For \textbf{\texttt{FileNotFoundError}}, check the working directory and filename. Fix the path rather than moving files into the project root.
\end{pitfall}

\section{Two ways to set up: plain Python or Anaconda}

\begin{keyidea}
Use one environment per project; activate it before installing or running packages.
\end{keyidea}

\section{Route A: install Python and configure PATH}

Download Python from the \href{https://www.python.org/downloads/}{\textbf{official downloads page}}.

\begin{description}[leftmargin=1.15in,style=nextline]
  \item[Windows] Run the installer, enable \textbf{\texttt{PATH}}, and reopen the terminal. Use \textbf{\texttt{py}}.
  \item[macOS] Run the \textbf{\texttt{.pkg} installer} and restart Terminal. Keep the \textbf{system Python} unchanged.
  \item[Linux] Install \textbf{Python 3} with your package manager. Do not replace the \textbf{system Python}.
\end{description}

\textbf{Windows PATH fix.} If \texttt{py} is missing, open \textbf{Edit environment variables for your account}, then \textbf{Path > Edit}. Add the installer's command directory, for example:

\begin{lstlisting}[style=mlterminal]
%LocalAppData%\Python\bin
\end{lstlisting}

\textbf{macOS PATH fix.} If the installer placed Python in \texttt{/usr/local/bin} but your terminal cannot find it, add that directory to your Z shell's \textbf{PATH}:

\begin{lstlisting}[style=mlterminal]
echo 'export PATH="/usr/local/bin:$PATH"' >> ~/.zshrc
source ~/.zshrc
\end{lstlisting}

For \textbf{Debian or Ubuntu}:

\begin{lstlisting}[style=mlterminal]
sudo apt update
sudo apt install python3 python3-venv python3-pip
\end{lstlisting}

\textbf{Verify} the installation and executable location:

\begin{lstlisting}[style=mlterminal]
# Windows
py --version
py -m pip --version
where py

# macOS or Linux
python3 --version
python3 -m pip --version
which -a python3
\end{lstlisting}

\section{Route A: create an isolated virtual environment}

\textbf{\texttt{venv}} keeps this project's packages separate from other projects. See the \href{https://docs.python.org/3/library/venv.html}{\textbf{venv documentation}}.

The root-level \textbf{\texttt{requirements.txt}} lists dependencies, including PyTorch (\texttt{torch}) for Labs~7--8. A CPU is sufficient.

\begin{lstlisting}[style=mlterminal]
numpy
pandas
matplotlib
scikit-learn
joblib
jupyter
torch
\end{lstlisting}

\textbf{Create} the environment in the \textbf{project root}:

\begin{lstlisting}[style=mlterminal]
# Windows
py -m venv .venv

# macOS or Linux
python3 -m venv .venv
\end{lstlisting}

\textbf{Activate} it from the same folder:

\begin{lstlisting}[style=mlterminal]
# Windows (PowerShell)
.venv\Scripts\Activate.ps1

# macOS/Linux
source .venv/bin/activate
\end{lstlisting}

\textbf{Install dependencies} after activation:

\begin{lstlisting}[style=mlterminal]
# Windows
py -m pip install --upgrade pip
py -m pip install -r requirements.txt

# macOS or Linux
python3 -m pip install --upgrade pip
python3 -m pip install -r requirements.txt
\end{lstlisting}

\textbf{Install once; activate every session.} Reuse \textbf{\texttt{.venv}} without reinstalling unchanged requirements.

\begin{mlfigure}
\begin{tikzpicture}[
  font=\small,
  pystep/.style={draw=MLBlue, fill=white, rounded corners=2pt, align=center,
               inner sep=4pt, minimum width=3.3cm, minimum height=1.05cm,
               font=\scriptsize\ttfamily\color{MLNavy}},
  pyonce/.style={pystep, fill=MLPaleBlue},
  pyevery/.style={pystep, fill=MLPaleTeal, draw=MLTeal},
  pycap/.style={font=\scriptsize\bfseries\color{MLMuted}},
  pynote/.style={font=\scriptsize\color{MLMuted}, align=center},
  pygo/.style={-{Stealth[length=2.2mm]}, thick, draw=MLTeal}]

\node[pyonce]  (venv) at (0,0)     {py -m venv .venv};
\node[pyevery] (act)  at (4.05,0)  {activate};
\node[pyonce]  (inst) at (8.10,0)  {pip install -r\\requirements.txt};
\node[pyevery] (work) at (12.15,0) {run notebooks in\\all\_experiments/};

\draw[pygo] (venv) -- (act);
\draw[pygo] (act)  -- (inst);
\draw[pygo] (inst) -- (work);

\node[pycap] at (0,0.92)     {once per project};
\node[pycap] at (4.05,0.92)  {every session};
\node[pycap] at (8.10,0.92)  {once};
\node[pycap] at (12.15,0.92) {every session};

% the return path: a later session rejoins at activation, not at install
\draw[pygo, draw=MLOrange]
  (work.south) -- ++(0,-0.85) -| (act.south);
\node[pynote, text=MLOrange] at (8.10,-1.15)
  {next session: \textbf{reactivate}, do not reinstall};

\end{tikzpicture}
\end{mlfigure}

\begin{keyidea}
\textbf{Activate first, then install.} Keep course packages inside the \textbf{project environment}, not in a shared global installation.
\end{keyidea}

\begin{mlfigure}
\begin{tikzpicture}[
  font=\small,
  envbox/.style={draw=MLBlue, fill=MLPaleBlue, rounded corners=2pt,
                 align=center, inner sep=4pt, minimum width=4.3cm},
  sysbox/.style={draw=MLMuted, fill=MLPaleGray, rounded corners=2pt,
                 align=center, inner sep=4pt, minimum width=4.3cm},
  pkg/.style={font=\scriptsize\ttfamily, align=center},
  lbl/.style={font=\scriptsize\color{MLMuted}, align=center},
  hit/.style={-{Stealth[length=2.2mm]}, thick, draw=MLTeal}]

\node[sysbox] (sys) at (0,2.45) {system Python};
\node[pkg,text=MLMuted,anchor=west] at (2.35,2.45) {leave this alone};

\node[envbox] (a) at (-2.80,0) {\texttt{THIS DIRECTORY/.venv}};
\node[pkg,anchor=north] at (-2.80,-0.52) {numpy\\pandas\\scikit-learn\\jupyter};

\node[envbox] (b) at (2.80,0) {\texttt{some other project/.venv}};
\node[pkg,anchor=north] at (2.80,-0.52) {its own versions\\of the same names};

\draw[hit] (sys) -- (a);
\draw[hit] (sys) -- (b);

\node[lbl,text width=\linewidth,inner xsep=0pt] at (0,-2.60)
  {\textbf{\color{MLNavy}Separate environments, separate packages.}
   Activate the intended \textbf{project environment} before running \texttt{pip install}.};

\end{tikzpicture}
\end{mlfigure}

\section{Route B: Anaconda or Miniconda}

\textbf{Miniconda} supplies Python and \texttt{conda}; \textbf{Anaconda} adds Navigator and preinstalled packages. Download \href{https://docs.anaconda.com/miniconda/}{\textbf{Miniconda}} or \href{https://www.anaconda.com/download}{\textbf{Anaconda}}.

\begin{description}[leftmargin=1.15in,style=nextline]
  \item[Windows] Use default installer options without adding conda to PATH. Open \textbf{Anaconda Prompt} from Start.
  \item[macOS / Linux] Run the \texttt{.pkg} or \texttt{.sh} installer, then close and reopen the terminal. The prompt shows \texttt{(base)}, which means conda is active.
\end{description}

\textbf{Create} one environment for the course, with a Python version the libraries support, then \textbf{activate} it:

\begin{lstlisting}[style=mlterminal]
conda create -n ml python=3.12
conda activate ml
\end{lstlisting}

\textbf{Install dependencies} with the same \texttt{requirements.txt} used in Route A:

\begin{lstlisting}[style=mlterminal]
pip install -r requirements.txt
\end{lstlisting}

Create once, activate each session, and avoid installing into \texttt{(base)}. Check the active environment and Python:

\begin{lstlisting}[style=mlterminal]
conda env list
python -c "import sys; print(sys.executable)"
\end{lstlisting}

\begin{pitfall}
Avoid installing the same package with both conda and pip. Use pip with the course requirements. Recreate a broken environment with \texttt{conda env remove -n ml}, then repeat setup.
\end{pitfall}

\begin{troubleshoot}
\textbf{Conda not found on Windows:} use Anaconda Prompt, or initialize PowerShell there with \texttt{conda init powershell}. \textbf{Jupyter import failure:} activate \texttt{ml}, restart Jupyter, and check the kernel with Lab~1's environment cell.
\end{troubleshoot}
